\scanpage{121}
\mspage{108}
\diagram{\begin{tikzcd}[ampersand replacement=\&,column sep=large]
(6.2.5.4) \& X(0)\arrow[r,dash,"{X(\unsure{d|\alpha})}"]\arrow[d,"{u(0)}"']\arrow[dr,phantom,"\text{cocart}"] \& \unsure{x|X}(1)\arrow[d]\arrow[dd,bend left=48,dash,dashed,"{u(1)}"]\\
\& Y(0)\arrow[r,"{\unsure{v|Y}}"]\arrow[dr,dashed,"{Y(\alpha)}"'] \& \widetilde Y\arrow[d,"\widetilde u"']\\
\& \& Y(1)
\end{tikzcd}}

Pour qu'on ait $u\in\Phi(\Delta^1)$, il f et s
que l'on ait
\[
 u(0),u(1),\widetilde u\ \in\ \Phi.\tag*{(6.2.5.5)}
\]

J'ai déjà vérifié la \unsure{chose} (sous forme
duale), mais ne retrouve plus l'endroit
(peut-être n'ai-je pas pris des notes).
Je n'ai pas envie de rechercher la
démonstration ici.

Pour achever de prouver 6.1. d),
i.e. l'égalité $C(I)=C_I$ pour $I$
artinien, j'ai envie de prouver encore
\unsure{qque} chose de plus général que voici :

\uline{Corollaire} 6.3. Soient $(M,\Phi,\Psi)$ \inserted{et $I\longmapsto\Phi(I)\unsure{\subset|\in}\operatorname{Fl}M(I)$} comme
dans 6.2.5. Soit $I$ un ens. ordonné,
\deleted{Alors \unsure{on}} \inserted{et} considérons aussi $\Phi_I$ défini (dans la
\deleted{(6.2.\unsure{1|2}} situation duale) \deleted{\unsure{page}} 25--(2.2) (mais
avec des notations duales, \unsure{i.e.}). Je dis que l'on
\[
 \Phi_I\supset\Phi(I),\tag*{(6.3.1)}
\]
% archive sheet number: 120

